% ID: SEL-D0012-BC-E05
% ESTADO: validada. S013; tres pasadas del mismo autor, sin segundo revisor.
\textbf{Datos y partición por plantas.}\par\nopagebreak[4]
Sea $D$ el suceso coche defectuoso. Las plantas forman una partición: $P(A)=0.45$, $P(B)=0.21$ y $P(C)=1-0.45-0.21=0.34$. Los complementos de las tasas de defecto son $0.99$, $0.97$ y $0.98$.

\textbf{a) No defectuoso y procedente de C.}\par\nopagebreak[4]
Se pide una intersección, por la regla del producto:
\[
\boxed{P(D^c\cap C)=P(C)P(D^c\mid C)=0.34\cdot0.98=0.3332.}
\]
\textbf{b) Planta A, condicionado a no ser defectuoso.}\par\nopagebreak[4]
Primero calculamos la probabilidad total de no defecto:
\[
\begin{aligned}
P(D^c)&=0.45\cdot0.99+0.21\cdot0.97+0.34\cdot0.98\\
&=0.4455+0.2037+0.3332=0.9824.
\end{aligned}
\]
Por Bayes,
\[
\boxed{P(A\mid D^c)=\frac{0.4455}{0.9824}
=\frac{4455}{9824}\approx0.453481.}
\]
Comprobación complementaria: $P(D)=0.0045+0.0063+0.0068=0.0176$ y $P(D)+P(D^c)=1$. No se intercambian $P(D^c\mid A)$ y $P(A\mid D^c)$.
